\section{Discussion}
\StableKG{} turns an answer's response to evidence withdrawal into an explicit inference property. The certificate captures how support is organized across temporal windows, a property that the unedited score vector alone cannot identify (Fig.~\ref{fig:interface}). Connecting this property to calibrated selection and locally maintained memory gives a common interface for prediction, acceptance and revision. The same evidence representation supports both the guarantee used to accept an answer and the computation used to revisit it.

The distinction between confidence and edit stability is central to interpreting the evidence. Confidence controls already rank correctness effectively for the stronger backbone. A certificate supplies a different capability: it identifies predictions whose answer is invariant to every allowed window deletion. The perturbation results show that the signed surplus also contains useful information about vulnerability beyond the original partition. Consequently, certification can be used as an explicit acceptance requirement even when its features add little to a correctness calibrator. This separation makes the system's decision interpretable: the probability concerns empirical correctness, and the certificate concerns a specified intervention.

This contribution complements recent temporal conformal prediction and interpretable graph editing. NCPNET addresses prediction-set coverage under non-exchangeable temporal observations \citep{wang2025nonexchangeable}; IGETR uses graph paths and language-model-guided editing to improve reasoning \citep{fan2026editing}. The present guarantee is a deterministic property of one fixed answer under memory withdrawal. Its exactness comes from a structured inference interface, whereas sampled robustness methods such as GraphProt and AdaptDel obtain guarantees through their own randomized constructions and perturbation models \citep{yang2025graphprot,huang2025adaptdel}. These differences suggest a modular division of labor: expressive encoders represent temporal knowledge, statistical selectors estimate answer reliability, and explicit evidence channels make selected revisions exactly analyzable.

The mixture also exposes a useful design choice. Increasing the contribution of temporal memory changes both predictive performance and sensitivity to that memory. The strong and compact backbones occupy different parts of this trade-off, so a common accuracy benefit is not required for a common certificate interface. Validation chooses the predictive operating point, while a deletion budget specifies the tolerated intervention. The guarantee applies to the frozen encoder, fixed admissible vocabulary and declared window family. Extending it to encoder retraining, arbitrary counter-evidence insertion or a continually changing candidate set requires additional analysis. Chronological forecasting and application-defined withdrawal windows are natural next tests of the interface.

Local maintenance makes this interpretation operational. The reverse index identifies exactly which registered queries depend on an edited event, and the measured comparison includes the writes and certificate refresh on both paths. The result is a system in which prediction, acceptance and revision share an explicit evidence representation. This provides a concrete basis for exposing neural knowledge-graph answers together with an inspectable tolerance to change.
